\documentclass[10pt,a4paper]{amsart}
\usepackage[margin=19mm]{geometry}
\usepackage{amsmath,amssymb,mathtools}
\usepackage[scr=rsfs,cal=euler]{mathalfa}
\usepackage{xfrac}
\usepackage[unicode,hidelinks,bookmarks=false]{hyperref}
\newcommand{\ie}{i.e.\ }
\newcommand{\cf}{cf.\ }
\newcommand{\resp}{respectively\ }
\newcommand{\Subsection}{\subsection}
\providecommand{\nobreakdash}{\nobreak\hbox{-}}
\setlength{\emergencystretch}{2em}
\multlinegap0pt
\usepackage{bm}
\usepackage{bbm}
\usepackage{dsfont}
\usepackage{xspace}
\usepackage{cancel}
\usepackage{mathtools}
\usepackage{amsthm}
\makeatletter
\@ifundefined{theorem}{\newtheorem{theorem}{Theorem}}{}
\@ifundefined{lemma}{\newtheorem{lemma}[theorem]{Lemma}}{}
\makeatother
\def\MM{\mathcal{M}}
\title[Optimal Harvesting in Stream Networks: Maximizing Biomass an]{Optimal Harvesting in Stream Networks: Maximizing Biomass and Yield}
\author{Tung D. Nguyen, Zhisheng Shuai, Tingting Tang, Amy Veprauskas, Yixiang Wu, Ying Zhou}
\date{Source: arXiv:2602.05071v1 (CC BY 4.0)}
\begin{document}
\maketitle

\noindent\emph{Source and licence.} Optimal Harvesting in Stream Networks: Maximizing Biomass and Yield. Version arXiv:2602.05071v1 (CC BY 4.0), distributed under the Creative Commons Attribution 4.0 International licence, as recorded in the arXiv licence metadata for this exact version and shown on the abstract page https://arxiv.org/abs/2602.05071v1. Full rights notice: licenses/arxiv-2602.05071-CC-BY-4.0.txt.\par\smallskip

\noindent\textit{Editorial scope.} Counted unit: the Theorem whose source label is \texttt{theorem:M'} in the article (source0.tex lines 360--375, source label \texttt{theorem:M'}), delivered together with the article's own complete original proof of it (source0.tex lines 377--423, one \texttt{proof} environment of 47 lines). No mathematics is written, repaired or re-derived: statement and proof are the article's own text line for line. Required context retained uncounted: eqn:ODEs\_2patch (lines 253--256); lem:persist (lines 258--264); eqn:u1 (lines 289--292); eqn:u2 (lines 289--292); lem:M'\_hetero (lines 299--311). Every definition, display and label of those lines is retained unchanged. Omitted from this package and not redistributed: 1--252, 257--257, 265--288, 293--298, 312--359, 376--376, 424--1148. Nothing inside a retained excerpt is omitted or altered: every retained body is delivered exactly as the article's lines are written, so no omission receipt of any kind accompanies this package. Citations: the retained excerpts carry no citation command at all, so no bibliography entry of the article is transcribed and no reference is redistributed. Reference adaptation: none was needed and none was made; every reference inside the delivered unit resolves inside it, so no unresolved reference or citation is produced. Printed slips kept literal: no printed word, symbol, hypothesis or number of the counted statement or of its proof was altered. Layout and class: the article's own class is \texttt{article} with packages including geometry, graphicx, amsmath, amsfonts, amssymb, amsthm, amstext, amsopn, appendix, caption, subcaption, tikz, booktabs, multirow; the wrapper is the lane's pinned \texttt{amsart} editorial wrapper at 10pt on A4, which restates the theorem-like environments the retained text needs and adds the article's own macro definitions that the retained text needs (\texttt{\textbackslash MM}), with extra wrapper packages (bm, bbm, dsfont, xspace, cancel, mathtools). The delivered numbering is the unit's own. Assets: no retained span contains a figure environment, a graphics-inclusion command, a PDF-page inclusion, a table or a TikZ picture; no image file is copied into this package, the package declares no asset dependency and no image file is redistributed with it. Licence: version arXiv:2602.05071v1 of this article is distributed under the Creative Commons Attribution 4.0 International licence, as recorded in the arXiv licence metadata for this exact version and shown on the abstract page https://arxiv.org/abs/2602.05071v1. A full rights notice is delivered at licenses/arxiv-2602.05071-CC-BY-4.0.txt. Characters: every retained excerpt is pure ASCII, so the delivered engine, XeLaTeX with its default Latin Modern text font, reports no missing character.
\begin{align}\label{eqn:ODEs_2patch}
    &u_1'=u_1(r_1-\theta H -c_1u_1) -(d+q)u_1+du_2\nonumber\\
    &u_2'=u_2(r_2-(1-\theta)H -c_2u_2)+(d+q)u_1-du_2.
\end{align}

\begin{lemma}\label{lem:persist}
    Suppose the intrinsic growth rates $r_1$ and $r_2$ satisfy
    \[
    \frac{d}{2d+q}r_1+\frac{d+q}{2d+q}r_2>H\frac{d+q}{2d+q} =: r_{\text{crit}}.
    \]
    Then there exists a unique  positive equilibrium $\bm u^*$ which is asymptotically stable.
\end{lemma}

\begin{align}
    &u_1(r_1-\theta H -c_1u_1) -(d+q)u_1+du_2=0\label{eqn:u1}\\
    &u_2(r_2-(1-\theta)H -c_2u_2)+(d+q)u_1-du_2=0\label{eqn:u2}.
\end{align}    

\begin{lemma}\label{lem:M'_hetero}
We have the equality
\begin{equation}\label{eq:M'_het}
    \MM'(\theta)= -\frac{H(u_1+u_2)((d+q)u_1^2-du_2^2)+Hu_1^2u_2^2(c_2-c_1)}{\det(A)},
\end{equation}
where 
\begin{equation}\label{matrix A_het}
A=\begin{bmatrix}
    -(c_1u_1^2+du_2) & du_1\\
    (d+q)u_2 & -(c_2u_2^2+(d+q)u_1)
\end{bmatrix}.
\end{equation}
\end{lemma}

\begin{theorem}\label{theorem:M'}
 Consider system \eqref{eqn:ODEs_2patch} and assume the persistence condition in Lemma \ref{lem:persist}. Assume further that $c_1=c_2=c$.
    \begin{enumerate}
        \item[(a)] Suppose the growth rates $r_1$ and $r_2$ satisfy 
        \[\frac{r_1\sqrt{d+q}-r_2\sqrt{d}}{\sqrt{d+q}-\sqrt{d}}> 2d+q +\frac{\sqrt{d+q}}{\sqrt{d+q}-\sqrt{d}}H=:r_M.\]
        Then $\MM'(\theta)< 0$ for any $\theta\in[0,1]$. Thus $\MM(\theta)$ is maximized at $\theta=0$, i.e. when the harvesting effort is concentrated on the downstream patch.
        \item[(b)]  Suppose the growth rates $r_1$ and $r_2$ satisfy 
        \[\frac{r_1\sqrt{d+q}-r_2\sqrt{d}}{\sqrt{d+q}-\sqrt{d}}< 2d+q -\frac{\sqrt{d}}{\sqrt{d+q}-\sqrt{d}}H=:r_m.\]
        Then $\MM'(\theta)> 0$ for any $\theta\in[0,1]$. Thus $\MM(\theta)$ is maximized at $\theta=1$, i.e. when the harvesting effort is concentrated on the upstream patch.
        \item[(c)] Suppose 
        \[
        r_m\leq \frac{r_1\sqrt{d+q}-r_2\sqrt{d}}{\sqrt{d+q}-\sqrt{d}}\leq r_M
        \]
        then there exists a unique $\theta^*\in[0,1]$ such that $\MM'(\theta^*)=0$. Furthermore, $\theta^*$ is a minimum of $\MM(\theta)$ in the interval $[0,1]$ and $\MM(\theta)$ is maximized at either $\theta=0$ or $\theta=1$, i.e. when the harvesting effort is concentrated  on either patch.
        \end{enumerate}
\end{theorem}

\begin{proof}
    Let $t:=\frac{u_1}{u_2}$. Dividing equations \eqref{eqn:u1} and \eqref{eqn:u2} by $u_1$ and $u_2$, respectively, yields
 \begin{subequations}
     \begin{align}
         &r_1-\theta H -(d+q) + \frac{d}{t}=cu_1 \label{eqn:div_u1_het}\\
         &r_2-(1-\theta)H -d + (d+q)t = cu_2. \label{eqn:div_u2_het}
     \end{align}
\end{subequations}
By taking \eqref{eqn:div_u1_het}$\div$\eqref{eqn:div_u2_het}, we obtain
\[
\frac{r_1-\theta H -(d+q) + \frac{d}{t}}{r_2-(1-\theta)H -d + (d+q)t}=t.
\]
Thus $t$ satisfies
\[
f(t,\theta) : = t(r_2-(1-\theta)H -d + (d+q)t) - (r_1-\theta H -(d+q) + \frac{d}{t}) =0.
\]
Consider a fixed value of $\theta$. Note that $r_2-(1-\theta)H -d + (d+q)t>0$ due to \eqref{eqn:div_u2_het}, so $f(t,\theta)$ is an increasing function with respect to $t$. 


From Lemma \ref{lem:M'_hetero} and the assumption that $c_1=c_2=c$, the sign of $\MM'(\theta)$ is the opposite of the sign of $(d+q)u_1^2-du_2^2=(d+q)u_2^2(t^2-\frac{d}{d+q})$. Thus we wish to examine the condition for $t>\sqrt{\frac{d}{d+q}}$ and $t<\sqrt{\frac{d}{d+q}}$. To this end, we use direct calculations and evaluate
\begin{equation}
    f\bigg(\sqrt{\frac{d}{d+q}},\theta\bigg) = \frac{\sqrt{d+q}-\sqrt{d}}{\sqrt{d+q}}\bigg(g(\theta)-\frac{r_1\sqrt{d+q}-r_2\sqrt{d}}{\sqrt{d+q}-\sqrt{d}}\bigg),
\end{equation}
where $g(\theta):=2d+q+\frac{\theta\sqrt{d+q}-(1-\theta)\sqrt{d}}{\sqrt{d+q}-\sqrt{d}}H$. It is easy to see that $g(\theta)$ is increasing in $\theta$. Furthermore, $g(0)=r_m$ and $g(1)=r_M$, where $r_m$ and $r_M$ are defined in Theorem \ref{theorem:M'}. We now consider three cases. 

\vspace{.1in}
\noindent\textbf{Case (a):} Suppose that
$\frac{r_1\sqrt{d+q}-r_2\sqrt{d}}{\sqrt{d+q}-\sqrt{d}}>r_M$.

Since $g(\theta)$ is increasing in $\theta$, we have $\frac{r_1\sqrt{d+q}-r_2\sqrt{d}}{\sqrt{d+q}-\sqrt{d}}>r_M=g(1)\geq g(\theta)$. As a result, 
$f\big(\sqrt{\frac{d}{d+q}},\theta\big) <0$, which implies $t>\sqrt{\frac{d}{d+q}}$. From Lemma \ref{lem:M'_hetero} and the assumption that $c_1=c_2=c$, we  conclude that $\MM'(\theta)<0$ and $\MM(\theta)$ is maximized at $\theta=0$.

\vspace{.1in}
\noindent\textbf{Case (b):} Suppose that
$\frac{r_1\sqrt{d+q}-r_2\sqrt{d}}{\sqrt{d+q}-\sqrt{d}}<r_m$.

Since $g(\theta)$ is increasing in $\theta$, we have $\frac{r_1\sqrt{d+q}-r_2\sqrt{d}}{\sqrt{d+q}-\sqrt{d}}<r_m=g(0)\leq g(\theta)$. As a result, 
$f\big(\sqrt{\frac{d}{d+q}},\theta\big) >0$, which implies $t<\sqrt{\frac{d}{d+q}}$. From Lemma \ref{lem:M'_hetero} and the assumption that $c_1=c_2=c$, we  conclude that $\MM'(\theta)>0$ and $\MM(\theta)$ is maximized at $\theta=1$.

\vspace{.1in}
\noindent\textbf{Case (c):} Suppose that $r_m\leq \frac{r_1\sqrt{d+q}-r_2\sqrt{d}}{\sqrt{d+q}-\sqrt{d}}\leq r_M$.
\vspace{.1in}

In this case, there exists a unique $\theta^*\in[0,1]$ such that $\frac{r_1\sqrt{d+q}-r_2\sqrt{d}}{\sqrt{d+q}-\sqrt{d}}=g(\theta^*)$. Thus $f\big(\sqrt{\frac{d}{d+q}},\theta^*\big) =0$, which  implies $t=\sqrt{\frac{d}{d+q}}$ when $\theta=\theta^*$. From Lemma \ref{lem:M'_hetero} and the assumption that $c_1=c_2=c$ we must have $\MM'(\theta^*)=0$. 

Furthermore, for any $\theta\in[0,\theta^*]$, we have $\frac{r_1\sqrt{d+q}-r_2\sqrt{d}}{\sqrt{d+q}-\sqrt{d}}>g(\theta)$ and thus $\MM'(\theta)<0$. Using an identical argument, we also have $\MM'(\theta)>0$ for any $\theta\in[\theta^*,1]$. Thus in the interval $[0,1]$, the total biomass $\MM(\theta)$ has only one local extremum $\theta^*$, which is a minimum. This also implies that $\MM(\theta)$ is maximized at either $\theta=0$ or $\theta=1$.
\end{proof}

\end{document}
